% ============================================================
% PAPER 7: Synonymix -- Privacy-Preserving Life Story Personas
% via Graph-Based Abstraction (Chen \& Kumar)
% ============================================================

% PERSONA_CATEGORY: Surrogate user persona (LLM-simulated participant)
% PERSONA_SUBCATEGORY: Persona-as-narrative-life-story profile (high-fidelity behavioural proxy for a real individual)
% PERSONA_SUBCATEGORY: Persona-as-privacy-preserving synthetic composite (graph-sampled recombination of many real people)
% PERSONA_SUBCATEGORY: Persona-as-collective artifact ("unigraph" for meso-level sensemaking)

\subsection*{Paper 7: \textsc{Synonymix} (Chen \& Kumar)}

This paper inherits the classical HCI framing of personas as ``a bridge
between user research and design decisions,'' but immediately relocates the
term into the generative-agent / agent-bank tradition: a persona is the
\emph{data grounding} that conditions an LLM agent so that it can act as a
behavioural proxy for a person. The authors' explicit contrast is between
\emph{demographic-only personas} --- criticised as ``flattened stereotypes''
that collapse within-group diversity into a single behavioural mode --- and
\emph{life story personas}, derived from McAdams-style narrative-psychology
life story interviews, which capture ``not just what happened to someone but
what it meant to them.'' Persona here is therefore not a design communication
device but a \emph{fidelity parameter}: the richer the personal narrative, the
better the simulated respondent.

The operational unit is unusual and is the paper's main contribution. Each
persona is first re-represented as a typed knowledge graph with Subject nodes
(people, places, institutions), Factual nodes (events, milestones) and
Interpretive nodes (values, motivations, self-narratives generated by LLM
``expert reflection'' in the voice of a sociologist, anthropologist or
psychologist). Individual persona graphs are then genericised (``Graduated from
Harvard University'' $\rightarrow$ ``Graduated from University'' $+$ a separate
subject node), semantically deduplicated with provenance tags, and merged into
a single collective structure, the \emph{unigraph}. Novel personas are then
\emph{sampled} from this structure by a ``Thematic Random Walk'' anchored on an
interpretive node --- producing ``synthetic individuals whose life stories are
composites of many real people but copies of none.'' Persona thus becomes a
traversal path through a shared experiential graph rather than a document, a
profile, or a demographic vector.

Evaluation reinforces this survey-respondent reading: three ``agent banks''
($N=30$ each: Demographic, Life Story, Synonymix) are seeded from General
Social Survey 2022 demographics and answer 108 non-demographic GSS items;
persona quality is measured as \emph{group-level distributional} fidelity
(EMD/TVD between response distributions) rather than as design usefulness or
empathy. Privacy is likewise operationalised at the persona level via Maximum
Source Contribution (the share of a synthetic persona's nodes traceable to its
most dominant real source, $<13\%$), with differential privacy over node
histograms proposed as future work.

Finally, the paper pushes the concept beyond the individual: the unigraph is
proposed as an interactive artifact for \emph{`meso-level'} interaction --- a
scale between a single AI persona and population-level modelling --- supporting
cross-temporal sensemaking (e.g., how ``the Class of 1985'' experienced shared
moments) and consensus building. Persona here shades into a \emph{queryable
collective representation}: not one representative user, but the structural
overlap of many real biographies. The authors also flag that privacy
preservation is ``necessary but not sufficient,'' leaving consent and community
agency in representation open, and warn that sparse graph coverage for
underrepresented groups may yield incoherent sampled life stories.

% ============================================================
% UPDATED RUNNING TAXONOMY TABLE
% ============================================================

\begin{table*}[t]
\centering
\caption{Running taxonomy of persona conceptualizations across workshop papers (updated through Paper 7).}
\label{tab:persona-taxonomy}
\small
\begin{tabular}{p{0.5cm} p{3.6cm} p{3.4cm} p{3.6cm} p{3.6cm}}
\toprule
\textbf{\#} & \textbf{Paper} & \textbf{Main category} & \textbf{Subcategory} & \textbf{Operationalization / unit} \\
\midrule
1 & AI Personas for UX Research in Large Scale Retail (Sharma et al., Walmart) &
Surrogate user persona (LLM-simulated participant) &
(a) Persona-as-test-participant; (b) Persona-as-latent-trait-vector (activation steering) &
Demographic/behavioural archetype profiles prompted into a multimodal LLM to produce synthetic usability feedback; validated against 150 human participants. Secondary: persona as steering direction in BLIP-2 activations \\
\addlinespace
2 & AI Personas in Highly Regulated Environments: Tools for Calibrating User Reliance? (Vasiliu \& Guarese) &
System-facing agent persona (deployed AI character) &
Persona-as-interaction-metaphor for trust/reliance calibration (incl.\ adaptive / metaphor-fluid personas) &
Two prompt-encoded conversational personas (\textit{Expert Operator} vs.\ \textit{Machine}) defining tone, stance and dialogue structure of a voice CAI; manipulated as an experimental variable and evaluated via trust/workload/usability measures and a behavioural disclosure probe \\
\addlinespace
3 & Evaluating Conversational Agent Integration in Peer Support Groups (Uetova et al.) &
Surrogate user persona (LLM-simulated participant) &
Persona-as-data-derived individual profile (behavioural clone) for multi-agent social simulation; explicitly contrasted with large ``persona catalogues'' (NLP persona datasets) &
One persona per real past participant: LLM-generated summary of that person's chat history, demographics and survey answers (tone, emoji use, thematic focus), used to condition synthetic peer replies in a replayed group chat; combined with probabilistic responder selection and ``buddy'' pairing; framed as a pre-deployment testbed for agent prompt/logic refinement, with stated limits on fidelity and inclusivity \\
\addlinespace
4 & Grounding Persona-Driven Usability Simulation in Established Standards (Touir, Barika Ktata \& Soui) &
Surrogate user persona (LLM-simulated participant) &
Persona-as-autonomous-task-agent (embodied evaluator in a perceive--decide--act loop); Persona-as-edge-case/accessibility probe &
Supervisor-Agent auto-synthesises personas from UI purpose + source code using a constrained schema (goal, tech-literacy level, accessibility constraint); each persona conditions a VLM/LLM User-Agent that clicks, scrolls and types on a live prototype, emitting action traces, error logs and think-aloud self-reports mapped to ISO 9241-110, Nielsen's heuristics and WCAG. Named exemplars (``Martha R.,'' ``David L.'') act as edge-case probes; governance layer forbids unsupported demographic stereotyping and requires provenance and human escalation \\
\addlinespace
5 & Ideal Persona AI in Real-time (Pham Thi Ngoc Anh) &
Self-expression persona (user's idealized self rendered by GenAI) &
(a) Persona-as-generated-self-image (real-time text-to-image self-portrait); (b) Persona-as-ethical-behaviour-envelope (transparent ``default ethical persona'' of the system) &
Persona defined as ``a personality belonging to a participant'' --- an ideal self co-authored with GenAI. Six designers wrote ideal-persona descriptions with ChatGPT, then fed five keywords/questions as prompts to StreamDiffusion in TouchDesigner and watched a live image of that persona; screen capture, field notes, think-aloud and post-study questions analysed thematically. Breakdowns: inconsistent persona prototypes, uncanny valley, disconnection/trust erosion. Proposes moving persona design from ``creative writing'' to ``structural engineering'': prompt training, declared default ethical persona, ethical constraints exposed as UI/UX controls, continuous user control in real time \\
\addlinespace
6 & Mapping Personas to Text Transformations: A Taxonomy Outline for Content Adaptation (Sillano, De Russis, Cal\`o, Troncy \& Lisena) &
Content-adaptation audience persona (NLP prompt-conditioning target reader) &
(a) Persona-as-decomposable-trait-vector mapped to auditable text-transformation operators; (b) Persona-as-target-reader-profile for content personalization &
Persona = specification of the intended reader that conditions an LLM rewrite, critiqued as a ``black-box prompt modifier''. Decomposed into four textually-correlated dimensions (demographic, domain knowledge, cognitive, contextual) synthesised from role-playing/personalization surveys; each dimension mapped in a taxonomy table to named operators (lexical, structural, content, style, tone) with worked examples. Two tabular persona profiles (P1 novice student, P2 science journalist) drive a five-step operator chain rewriting a thermodynamics sentence with GPT-5.4, each edit traceable to an operator + persona dimension. Goal: transparency, bias auditability and writer agency; mappings not yet empirically validated \\
\addlinespace
7 & \textsc{Synonymix}: Privacy-Preserving Life Story Personas via Graph-Based Abstraction (Chen \& Kumar) &
Surrogate user persona (LLM-simulated participant) &
(a) Persona-as-narrative-life-story profile (high-fidelity behavioural proxy for a real individual); (b) Persona-as-privacy-preserving synthetic composite (graph-sampled from many real people); (c) Persona-as-collective artifact (``unigraph'' for meso-level sensemaking) &
Persona framed as data grounding for a generative agent, on a fidelity continuum from ``flattened'' demographic personas to McAdams-style life story personas. Each persona is converted into a typed knowledge graph (Subject / Factual / Interpretive nodes, with I-nodes produced by LLM ``expert reflection''), genericised and semantically deduplicated into a merged \emph{unigraph}; new personas are sampled by a Thematic Random Walk anchored on an interpretive node. Evaluated as three 30-agent banks (Demographic, Life Story, Synonymix) answering 108 GSS 2022 items, scored by EMD/TVD distributional distance ($p<0.001$, $r=0.59$ on ordinal items) and by Maximum Source Contribution ($\overline{MSC}=0.129$) as a structural privacy metric; DP over node histograms deferred. Persona additionally reframed as a shareable collective representation supporting cross-temporal sensemaking and consensus building rather than prediction \\
\bottomrule
\end{tabular}
\end{table*}
